\subsection{Reached-set quiescence}

\begin{theorem}[Collatz settlement realised ($Cl = id$)]
  \label{thm:collatz-settlement-realised}
  \lean{Obsmetrologia.Closure.Collatz.settling_accumulates}
  \leanok
  \uses{thm:valuation, thm:settle-step, thm:closure-rule-lfp}
  The Collatz settling $F_U(E)_\tau = Cl(E)\vee_k u_\tau$ with $Cl=id$ is not a
  new tier: Obsmetrologia's \texttt{Valuation} already hosts the infinite entity set
  (\texttt{Entity = Mark} is string-backed), and the increment-driven flow
  \texttt{ModalFrame.drive} is $F_U$. The discrete frame
  (fiber $\varphi$ = [$\varphi$]) realizes $Cl=id$ (\texttt{closureOp\_discrete}),
  so the Collatz settling is pure reached-cell accumulation:
  \texttt{settling (n+1) = settling n $\vee_k$ u_n}.
\end{theorem}

\begin{theorem}[Stopping time is the closure ordinal under descent]
  \label{thm:collatz-kappa-stopping-time}
  \lean{Obsmetrologia.Closure.Collatz.reachedQuiesces_iff_reachesTerminal_of_all_fall_below}
  \leanok
  \uses{thm:collatz-settlement-realised, thm:collatz-quiesce-of-terminal, thm:collatz-reduction}
  The Collatz reached-set --- realized over Obsmetrologia's internal set theory
  \texttt{Discrete.Set}~$\mathbb{N}$ --- quiesces at a finite stage iff the seed
  reaches the terminal cycle: the closure ordinal $\kappa$ is the total stopping
  time under the universal descent hypothesis. Both directions are proved, and
  $\mathsf{Classical}$-free
  (revisit uses propositional equality \texttt{Nat.decEq}, not Bool
  \texttt{==}). The converse goes through the descent half
  (\cref{thm:collatz-reduction}): quiescence forces an orbit revisit, hence
  finitely many values with a global minimum $V$ that cannot fall below itself,
  so $V\le 2$.
\end{theorem}

\begin{theorem}[Closure ordinal stage is total stopping stage]
  \label{thm:collatz-kappa-is-total-stopping-stage}
  \lean{Obsmetrologia.Closure.Collatz.closureOrdinalStage_iff_totalStoppingStage}
  \leanok
  \uses{thm:collatz-settlement-realised}
  The finite closure ordinal stage of the reached-set settlement --- the first
  \texttt{ReachedQuiesces} stage --- is exactly the total stopping stage in the
  settlement sense: the first orbit value already seen. This theorem is proved without
  the descent frontier; it identifies the stage counted by $\kappa$ with the local
  stopping/revisit stage of the Collatz settlement.
\end{theorem}

\begin{theorem}[Finite closure ordinal stage exists exactly at terminal reachability]
  \label{thm:collatz-kappa-stage-exists}
  \lean{Obsmetrologia.Closure.Collatz.closureOrdinalStage_exists_iff_reachesTerminal_of_all_fall_below}
  \leanok
  \uses{thm:collatz-kappa-is-total-stopping-stage, thm:collatz-kappa-stopping-time}
  Under universal descent, existence of a finite \texttt{ClosureOrdinalStage} is
  equivalent to terminal reachability. Lean first proves
  \texttt{exists\_closureOrdinalStage\_iff\_reachedQuiesces}: an arbitrary
  quiescence witness yields a least orbit-revisit stage, avoiding any decidable
  equality assumption for the infinite reached-set functions.
\end{theorem}

\begin{theorem}[Total stopping stage exists exactly at terminal reachability]
  \label{thm:collatz-total-stopping-stage-exists}
  \lean{Obsmetrologia.Closure.Collatz.totalStoppingStage_exists_iff_reachesTerminal_of_all_fall_below}
  \leanok
  \uses{thm:collatz-kappa-is-total-stopping-stage, thm:collatz-kappa-stage-exists}
  Under universal descent, existence of a finite \texttt{TotalStoppingStage} is
  equivalent to terminal reachability. This is the existence-level total stopping
  time form of \(\kappa\): \texttt{exists\_totalStoppingStage\_iff\_closureOrdinalStage}
  transports the finite closure-ordinal stage across the stagewise equivalence.
\end{theorem}

\begin{theorem}[Residue-gap stopping-time bridge]
  \label{thm:collatz-residue-kappa-stopping-time}
  \lean{Obsmetrologia.Closure.Collatz.reachedQuiesces_iff_reachesTerminal_of_residueGap_and_crossing}
  \leanok
  \uses{thm:collatz-kappa-stopping-time, thm:collatz-residue-conditional}
  The residue inequality and crossing hypothesis identify reached-set quiescence with terminal
  reachability. This is the settlement-side face of the proved conditional route: once the
  descent skeleton is supplied by those two inputs, $\kappa$ is exactly the finite stopping
  stage of the reached-set settlement.
\end{theorem}

\begin{theorem}[Residue gap plus no-infinite-stream stopping-time bridge]
  \label{thm:collatz-residue-no-infinite-stream-kappa}
  \lean{Obsmetrologia.Closure.Collatz.reachedQuiesces_iff_reachesTerminal_of_residueGap_and_noPositiveInfiniteBarrierStream}
  \leanok
  \uses{thm:collatz-residue-kappa-stopping-time, thm:collatz-crossing-no-infinite-stream}
  The residue inequality and the no-positive-infinite-barrier-stream condition identify
  reached-set quiescence with terminal reachability. This is the kappa/stopping-time bridge
  stated in the same two-frontier language as \cref{thm:collatz-residue-no-infinite-stream}.
\end{theorem}

\begin{theorem}[Residue gap plus no-infinite-stream closure-ordinal stage]
  \label{thm:collatz-residue-no-infinite-stream-kappa-stage}
  \lean{Obsmetrologia.Closure.Collatz.closureOrdinalStage_exists_iff_reachesTerminal_of_residueGap_and_noPositiveInfiniteBarrierStream}
  \leanok
  \uses{thm:collatz-kappa-stage-exists, thm:collatz-residue-no-infinite-stream-kappa}
  In the two-frontier arithmetic language, the residue inequality plus absence of
  positive infinite barrier streams identifies terminal reachability with existence
  of a finite closure-ordinal stage. This is the explicit \(\kappa\)-stage form of
  the settlement bridge.
\end{theorem}

\begin{theorem}[Packaged frontier closure-ordinal stage]
  \label{thm:collatz-frontier-inputs-kappa-stage}
  \lean{Obsmetrologia.Closure.Collatz.closureOrdinalStage_exists_iff_reachesTerminal_of_frontierInputs}
  \leanok
  \uses{thm:collatz-frontier-inputs, thm:collatz-kappa-stage-exists}
  The same \(\kappa\)-stage bridge stated over \texttt{CollatzFrontierInputs}: once
  the modular residue gap and no-positive-infinite-barrier-stream inputs are supplied,
  terminal reachability is equivalent to existence of a finite closure-ordinal stage.
\end{theorem}

\begin{theorem}[Packaged frontier total stopping stage]
  \label{thm:collatz-frontier-inputs-total-stopping-stage}
  \lean{Obsmetrologia.Closure.Collatz.totalStoppingStage_exists_iff_reachesTerminal_of_frontierInputs}
  \leanok
  \uses{thm:collatz-frontier-inputs-kappa-stage, thm:collatz-total-stopping-stage-exists}
  Over \texttt{CollatzFrontierInputs}, terminal reachability is equivalent to
  existence of a finite total stopping stage. This is the packaged frontier form
  of the statement that the total stopping time is the closure ordinal \(\kappa\).
\end{theorem}

\begin{theorem}[One-frontier kappa package criterion]
  \label{thm:collatz-one-frontier-kappa-package}
  \lean{Obsmetrologia.Closure.Collatz.residueGap_alone_solves_kappa_package_iff_residueGap_implies_noStream}
  \leanok
  \uses{thm:collatz-one-frontier-criterion, thm:collatz-frontier-inputs-total-stopping-stage}
  A one-frontier reduction to the modular residue gap gives the whole
  terminal/\(\kappa\)/total-stopping package exactly when that same residue gap
  implies \texttt{NoPositiveInfiniteBarrierStream}. This theorem is the precise
  settlement-level test for the requested collapse to the single inequality
  \(C < D\cdot r\): the collapse holds only once the no-stream
  consequence is proved from the modular lock.
\end{theorem}
\begin{theorem}[Parent-gap kappa package criterion]
  \label{thm:collatz-parent-gap-kappa-package}
  \lean{Obsmetrologia.Closure.Collatz.parentBarrierGap_alone_solves_kappa_package_iff_parentBarrierGap_implies_noStream}
  \leanok
  \uses{thm:collatz-one-frontier-kappa-package, thm:collatz-parent-gap-one-frontier-criterion}
  At the original semigroup face, \texttt{ParentBarrierGap} alone gives the whole
  terminal/\(\kappa\)/total-stopping package exactly when it rules out positive
  infinite barrier streams.
\end{theorem}
\begin{theorem}[Named one-lock kappa package criterion]
  \label{thm:collatz-parent-one-lock-kappa-package}
  \lean{Obsmetrologia.Closure.Collatz.parentBarrierOneLockCollapse_iff_parentBarrierGap_alone_solves_kappa_package}
  \leanok
  \uses{thm:collatz-parent-gap-kappa-package, thm:collatz-parent-one-lock-collapse}
  The named collapse proposition \texttt{ParentBarrierOneLockCollapse} is exactly
  the statement that \texttt{ParentBarrierGap} alone gives terminal reachability
  together with finite closure-ordinal and total-stopping stages.
\end{theorem}
\begin{theorem}[Named one-lock stopping-time bridge]
  \label{thm:collatz-parent-one-lock-kappa-bridge}
  \lean{Obsmetrologia.Closure.Collatz.reachedQuiesces_iff_reachesTerminal_of_parentBarrierGap_and_oneLock}
  \lean{Obsmetrologia.Closure.Collatz.closureOrdinalStage_exists_iff_reachesTerminal_of_parentBarrierGap_and_oneLock}
  \lean{Obsmetrologia.Closure.Collatz.totalStoppingStage_exists_iff_reachesTerminal_of_parentBarrierGap_and_oneLock}
  \leanok
  \uses{thm:collatz-parent-one-lock-supplies-descent, thm:collatz-kappa-stopping-time, thm:collatz-kappa-stage-exists, thm:collatz-total-stopping-stage-exists}
  Given the semigroup gap and the named one-lock proposition, reached-set
  quiescence, finite closure-ordinal-stage existence, and finite total-stopping
  stage existence are each equivalent to terminal reachability. This is the
  direct \(\kappa\)/total-stopping bridge for the one-lock route.
\end{theorem}
\begin{theorem}[Named one-lock kappa package, forward form]
  \label{thm:collatz-parent-one-lock-kappa-package-forward}
  \lean{Obsmetrologia.Closure.Collatz.kappaPackage_of_parentBarrierGap_and_oneLock}
  \lean{Obsmetrologia.Closure.Collatz.kappaPackage_of_residueGap_and_oneLock}
  \leanok
  \uses{thm:collatz-parent-one-lock-kappa-bridge, thm:collatz-parent-one-lock-kappa-package, thm:collatz-residue-gap-reduces}
  Forward use of the named one-lock collapse: from the parent-barrier gap and
  \texttt{ParentBarrierOneLockCollapse}, Lean gets terminal reachability together
  with finite \texttt{ClosureOrdinalStage} and \texttt{TotalStoppingStage}
  witnesses for every positive seed. The residue version is the same statement
  over the modular inequality \(C<D\cdot r\).
\end{theorem}

\begin{theorem}[Residue-frontier closure-ordinal stage]
  \label{thm:collatz-residue-frontier-kappa-stage}
  \notready
  \uses{thm:collatz-residue-gap-holds, thm:collatz-crossing-holds, thm:collatz-kappa-stage-exists}
  The frontier version of the \(\kappa\)-stage bridge: the modular
  residue-gap frontier plus the crossing frontier identifies terminal reachability
  with existence of a finite \texttt{ClosureOrdinalStage}. It is red only through
  those frontier inputs.
\end{theorem}

\begin{theorem}[Unconditional stopping-time bridge]
  \label{thm:collatz-unconditional-kappa-stopping-time}
  \notready
  \uses{thm:collatz-kappa-stopping-time, thm:collatz-gap-holds, thm:collatz-crossing-holds}
  The unconditional version is the same bridge after substituting the two current frontier
  holes. It is red because it inherits \cref{thm:collatz-gap-holds} and
  \cref{thm:collatz-crossing-holds}; the conditional bridge above is the proved machinery.
\end{theorem}

\begin{theorem}[Revisit makes the reached-set quiesce]
  \label{thm:collatz-quiesce-of-terminal}
  \lean{Obsmetrologia.Closure.Collatz.quiesce_of_revisit}
  \leanok
  \uses{thm:collatz-settlement-realised}
  If the orbit revisits an already-reached value, the reached-set settles (adds
  nothing new); over \texttt{Discrete.Set}~$\mathbb{N}$ this follows directly from
  the Bool-union law. With the terminal cycle 2-periodic ($T^2(1)=1$,
  $T^2(2)=2$) it gives the easy half of the bridge: reaching $\{1,2\}$ forces
  quiescence.
\end{theorem}
